\documentclass[aspectratio=169]{beamer}
\input{header}
\coursecode{JKSSB}

\title{What Is a Computer?}
\subtitle{Lesson 01 --- Fundamentals \& Data Representation}
\author{JKSSB Computer Awareness}
\date{\today}

\begin{document}
\frame{\titlepage}

\begin{frame}{Video Lesson 01}
  \centering
  \Large\href{https://youtu.be/EdRZK57oqpo}{Watch: What is a Computer?}
\end{frame}

\begin{frame}{Start with a Familiar Task}
  \begin{itemize}
    \item A clerk enters marks into a spreadsheet.
    \item A formula totals the marks and calculates a percentage.
    \item The screen shows the result; the file can be saved for later.
  \end{itemize}
  \vspace{0.4cm}
  \begin{block}{Question}
    Which parts of this task are input, processing, output, and storage?
  \end{block}
\end{frame}

\begin{frame}{A Working Definition}
  \begin{block}{Computer}
    An electronic, programmable system that accepts data, processes it under
    instructions, produces information, and can store data or results.
  \end{block}
  \begin{itemize}
    \item \key{Data}: raw facts or symbols supplied to a system.
    \item \key{Information}: data organized or processed to be useful.
    \item A computer follows instructions; it does not guarantee that the
      supplied data or chosen instructions are correct.
  \end{itemize}
\end{frame}

\begin{frame}{Data Becomes Information}
  \begin{columns}[T]
    \column{0.46\textwidth}
      \textbf{Data}
      \begin{itemize}
        \item 68, 74, 81
        \item Three separate marks
        \item Not yet a summary
      \end{itemize}
    \column{0.46\textwidth}
      \textbf{Information}
      \begin{itemize}
        \item Total = 223
        \item Average = 74.33
        \item Results now answer a question
      \end{itemize}
  \end{columns}
  \vspace{0.3cm}
  \muted{The same values can serve as data in one task and information in
  another, depending on the question being asked.}
\end{frame}

\begin{frame}{The IPO Cycle}
  \begin{enumerate}
    \item \key{Input}: data and instructions enter through input devices or
      software.
    \item \key{Processing}: the CPU and programs transform the data.
    \item \key{Output}: results are presented through a display, printer,
      speaker, or another output device.
    \item \key{Storage}: data and results are kept for current or later use.
  \end{enumerate}
  \begin{alertblock}{Exam cue}
    Storage is often shown alongside the Input--Process--Output (IPO) model;
    the four-part view is commonly called IPOS.
  \end{alertblock}
\end{frame}

\begin{frame}{Characteristics: What Computers Do Well}
  \begin{itemize}
    \item \key{Speed}: perform many operations quickly.
    \item \key{Accuracy}: repeat calculations consistently when data and
      instructions are correct.
    \item \key{Diligence}: do not tire from repetitive work.
    \item \key{Storage}: retain large amounts of data in a compact form.
    \item \key{Versatility}: support different tasks by running different
      programs.
    \item \key{Automation}: continue a programmed task once started and
      supplied with required input.
  \end{itemize}
\end{frame}

\begin{frame}{Limitations: What Computers Do Not Decide}
  \begin{itemize}
    \item A computer does not independently determine whether an input is
      truthful or fair.
    \item Incorrect data or instructions can produce incorrect results.
    \item It has no human judgement or common-sense understanding by default.
    \item Output still needs interpretation by a person in context.
  \end{itemize}
  \begin{block}{Remember}
    \key{GIGO} means \emph{Garbage In, Garbage Out}: poor input can lead to
    poor output.
  \end{block}
\end{frame}

\begin{frame}{Applications Around Us}
  \begin{itemize}
    \item \key{Office}: documents, calculations, records, and communication.
    \item \key{Banking}: account records, transaction processing, and ATMs.
    \item \key{Education}: learning platforms, assessments, and digital
      resources.
    \item \key{Healthcare}: patient records, appointments, and diagnostic
      support.
    \item \key{Governance}: online forms, records, and citizen services.
  \end{itemize}
  \vspace{0.3cm}
  \muted{The application changes; the input--processing--output idea remains.}
\end{frame}

\begin{frame}{Quick Classification Drill}
  \begin{itemize}
    \item Raw marks entered into a form: \key{data / input}.
    \item A formula calculating an average: \key{processing}.
    \item The average shown on screen: \key{output / information}.
    \item A saved spreadsheet: \key{storage}.
  \end{itemize}
  \begin{exampleblock}{Takeaway}
    Identify the role in the current task. Do not classify a file as
    processing, or a screen as storage, merely because both are used by a
    computer.
  \end{exampleblock}
\end{frame}

\begin{frame}{Lesson 01: Recall}
  \begin{itemize}
    \item A computer accepts data, follows instructions, processes, outputs,
      and can store.
    \item Data are raw facts; information is useful processed or organized
      data.
    \item IPO = Input, Processing, Output; storage is the additional S in
      IPOS.
    \item Speed and accuracy are strengths; correctness depends on input and
      instructions.
  \end{itemize}
\end{frame}
\end{document}
